\documentclass[11pt,a4paper]{article}
\usepackage[T1]{fontenc}\usepackage{mathptmx}\usepackage[a4paper,margin=20mm]{geometry}
\usepackage{graphicx,booktabs,array,tabularx,float,xcolor,enumitem}\usepackage[font=small,labelfont=bf]{caption}
\setlist{itemsep=1pt,topsep=2pt}\setlength{\parskip}{4pt}\setlength{\parindent}{0pt}
\newcolumntype{Y}{>{\raggedright\arraybackslash}X}
\title{\vspace{-12mm}Post-hoc analyses of the cross-validated copy-paste experiment\\\large Error analysis, detection versus segmentation, and augmentation versus model capacity}
\author{EGH490 -- analysis notes (12 October 2026)}\date{}
\begin{document}\maketitle\vspace{-8mm}

\section*{0. Data, validity and common methods}
\begin{itemize}
\item \textbf{No retraining or new inference.} All analyses use the test-fold predictions saved by the 120 finished models
(6 set-ups $\times$ 2 arms $\times$ 5 folds $\times$ 2 seeds): every box with confidence $\geq 0.001$ (top-1 class, NMS IoU 0.7,
$\le$300 per frame), and for segmentation the mask IoU of every prediction with every labelled fish. A CPU-only job re-packed these
files, the ground truth, the training logs and the model sizes. The saved predictions contain only each box's top class, so
second-choice species cannot be analysed.
\item \textbf{Same partitions, no leakage.} Each frame is scored only by the model whose test fold contains it (checked for every
prediction), using \texttt{folds.csv} and the ground-truth files of \texttt{cv\_evaluate.py}. Recomputing the published values from
the re-packed predictions reproduced them exactly (e.g.\ medium detection focus-species AP50 0.335/0.424; medium segmentation mask
AP50 0.308/0.458). A conservative check found that 7 of the 203 focus fish (all Orangespine Unicornfish, in two long encounters)
have other frames of the same encounter in the training/validation pool of the model that scores them; all other focus fish are
from encounters entirely outside that pool.
\item \textbf{Uncertainty.} Paired block bootstrap: the 49 blocks of 30 frames resampled 2,000 times with the generator and seed used
by \texttt{cv\_evaluate.py}; the same resample is applied to every model, so baseline/copy-paste, detection/segmentation and
small/medium/large comparisons are all paired by block. Seeds are averaged within each resample. Intervals are 95\% percentile
intervals; $p$ is the two-sided bootstrap $p$. Seed-to-seed variance is \emph{not} included.
\item \textbf{Status.} These analyses were \emph{not} pre-registered. They are exploratory and the intervals are not adjusted for
multiple comparisons; they should be reported as secondary and descriptive.
\end{itemize}

\section{Error analysis of the three focus species}
\textbf{Operating point and definitions.} Confidence threshold $\tau = 0.25$ (the threshold of the reported precision and recall);
box IoU threshold 0.5 for all set-ups (so detection and segmentation are compared on the same criterion). Predictions with
confidence $\geq\tau$ are matched to labelled fish of the same class, highest confidence first (as in COCO/\texttt{cv\_evaluate.py}).
Each labelled focus fish then falls into exactly one category, checked in this order:
\begin{enumerate}[label=(\alph*)]
\item \emph{Correct species} (TP): matched as above.
\item \emph{Wrong species} (classification error): an unmatched prediction of another species with IoU $\geq 0.5$ and confidence
$\geq\tau$ (predictions already matched to a neighbouring fish are excluded).
\item \emph{Poorly localised} (localisation error): an unmatched prediction with $0.1 \le$ IoU $<0.5$ and confidence $\geq\tau$.
\item \emph{Missed}: no prediction with IoU $\geq 0.1$ and confidence $\geq\tau$. Missed fish are further split by whether a
\emph{correct-species} box with IoU $\geq 0.5$ exists below $\tau$. These sub-threshold scores are reported only as evidence of a weak
signal, never as classifications (their median confidence is 0.02).
\end{enumerate}
False positives (unmatched focus-species predictions, confidence $\geq\tau$) are split into duplicates (IoU $\geq0.5$ with an
already-matched fish of that species), wrong-species FPs (IoU $\geq0.5$ with a labelled fish of another species), poorly localised
FPs, and FPs on \emph{unlabelled regions} (IoU $<0.1$ with every labelled fish). The last group may contain real but unlabelled
fish; it was not re-annotated and should not be read as hallucination. Results below pool all six set-ups and both seeds (12
models per arm; per-set-up results are in \texttt{components\_by\_setup.csv} and \texttt{err\_summary\_all.csv}).

\begin{table}[H]\centering
\caption{Fate of every labelled focus fish, baseline $\rightarrow$ copy-paste (pooled over 12 models per arm; 93, 77 and 33 fish).
Second line: change with 95\% block-bootstrap interval; bold: interval excludes zero. pp = percentage points.}
\label{tab:err}\footnotesize
\begin{tabularx}{\linewidth}{@{}p{4.4cm}YYY@{}}\toprule
Quantity & Forcepsfish & Orangespine Unicornfish & Convict Tang \\\midrule
Recall (\%) & 26.8 $\rightarrow$ 40.3 \newline \footnotesize{\textbf{$+$13.5 [$+$6.8, $+$23.2]}} & 15.5 $\rightarrow$ 19.6 \newline \footnotesize{\textbf{$+$4.1 [$+$1.6, $+$9.3]}} & 35.6 $\rightarrow$ 58.8 \newline \footnotesize{\textbf{$+$23.2 [$+$16.8, $+$34.9]}} \\
Precision (\%) & 57.7 $\rightarrow$ 60.3 \newline \footnotesize{$+$2.6 [$-$3.0, $+$10.8]} & 36.3 $\rightarrow$ 41.4 \newline \footnotesize{$+$5.1 [$-$0.3, $+$13.2]} & 59.0 $\rightarrow$ 53.3 \newline \footnotesize{$-$5.7 [$-$13.1, $+$7.1]} \\
Wrong species (\%) & 24.7 $\rightarrow$ 17.2 \newline \footnotesize{\textbf{$-$7.5 [$-$14.3, $-$2.1]}} & 15.4 $\rightarrow$ 14.2 \newline \footnotesize{$-$1.2 [$-$3.8, $+$3.1]} & 11.6 $\rightarrow$ 7.6 \newline \footnotesize{\textbf{$-$4.0 [$-$8.3, $-$1.5]}} \\
Poorly localised (\%) & 8.2 $\rightarrow$ 5.6 \newline \footnotesize{\textbf{$-$2.6 [$-$5.0, $-$0.4]}} & 2.7 $\rightarrow$ 1.8 \newline \footnotesize{$-$0.9 [$-$2.6, $+$0.2]} & 0.3 $\rightarrow$ 0.3 \newline \footnotesize{$+$0.0 [$+$0.0, $+$0.0]} \\
Missed, no prediction $\geq$0.25 (\%) & 40.2 $\rightarrow$ 36.8 \newline \footnotesize{$-$3.4 [$-$8.2, $+$0.2]} & 66.5 $\rightarrow$ 64.4 \newline \footnotesize{$-$2.1 [$-$10.5, $+$1.6]} & 52.5 $\rightarrow$ 33.3 \newline \footnotesize{\textbf{$-$19.2 [$-$28.3, $-$14.3]}} \\
\quad of which correct species $<$0.25 (\%) & 25.8 $\rightarrow$ 20.7 \newline \footnotesize{\textbf{$-$5.1 [$-$9.5, $-$1.1]}} & 37.1 $\rightarrow$ 31.9 \newline \footnotesize{\textbf{$-$5.2 [$-$13.5, $-$1.6]}} & 32.1 $\rightarrow$ 21.2 \newline \footnotesize{$-$10.9 [$-$25.9, $+$6.9]} \\
Localised, any species (\%) & 51.5 $\rightarrow$ 57.5 \newline \footnotesize{\textbf{$+$6.0 [$+$2.6, $+$10.8]}} & 30.8 $\rightarrow$ 33.8 \newline \footnotesize{$+$2.9 [$-$1.6, $+$10.8]} & 47.2 $\rightarrow$ 66.4 \newline \footnotesize{\textbf{$+$19.2 [$+$14.3, $+$28.3]}} \\
Correct species $|$ localised (\%) & 52.0 $\rightarrow$ 70.1 \newline \footnotesize{\textbf{$+$18.1 [$+$7.6, $+$33.8]}} & 50.2 $\rightarrow$ 58.0 \newline \footnotesize{$+$7.8 [$-$1.1, $+$14.2]} & 75.4 $\rightarrow$ 88.6 \newline \footnotesize{\textbf{$+$13.2 [$+$3.9, $+$37.2]}} \\
FP on other species (per model) & 6.6 $\rightarrow$ 11.2 \newline \footnotesize{\textbf{$+$4.7 [$+$1.0, $+$9.4]}} & 7.8 $\rightarrow$ 6.5 \newline \footnotesize{$-$1.2 [$-$3.3, $+$0.8]} & 4.0 $\rightarrow$ 9.6 \newline \footnotesize{\textbf{$+$5.6 [$+$1.2, $+$11.8]}} \\
FP on unlabelled regions (per model) & 7.8 $\rightarrow$ 9.4 \newline \footnotesize{$+$1.6 [$-$1.7, $+$4.7]} & 12.6 $\rightarrow$ 14.3 \newline \footnotesize{$+$1.8 [$-$1.8, $+$5.4]} & 3.2 $\rightarrow$ 6.9 \newline \footnotesize{\textbf{$+$3.7 [$+$1.5, $+$6.2]}} \\
All FP (per model) & 18.2 $\rightarrow$ 24.7 \newline \footnotesize{\textbf{$+$6.4 [$+$0.7, $+$13.0]}} & 20.9 $\rightarrow$ 21.3 \newline \footnotesize{$+$0.4 [$-$3.9, $+$4.4]} & 8.2 $\rightarrow$ 17.0 \newline \footnotesize{\textbf{$+$8.8 [$+$3.9, $+$15.4]}} \\
\bottomrule\end{tabularx}\end{table}

\begin{figure}[H]\centering\includegraphics[width=\linewidth]{figs/fig_err_fate.pdf}
\caption{What happened to every labelled focus fish, by species and arm (confidence $\geq0.25$, box IoU $\geq0.5$).}\label{fig:fate}\end{figure}

\begin{table}[H]\centering\caption{Most frequent species confusions (counts per model, mean of 12 models; confidence $\geq0.25$,
IoU $\geq0.5$). $^\dagger$Class spelt ``Lognnose'' in the label set.}\label{tab:conf}\footnotesize
\begin{tabularx}{\linewidth}{@{}p{3.3cm}p{3.6cm}YY@{}}\toprule
Species & Direction & Baseline & Copy-paste\\\midrule
Forcepsfish & Labelled as (FN) & Yellow Tang 22.7; Lined Butterflyfish 0.1 & Yellow Tang 15.6; Convict Tang 0.2 \\
 & Other fish labelled as it (FP) & Yellow Tang 3.2; Longnose Butterflyfish$^\dagger$ 1.8 & Yellow Tang 5.0; Longnose Butterflyfish$^\dagger$ 2.2 \\
Orangespine Unicornfish & Labelled as (FN) & Brown Surgeonfish 6.8; Goldring Surgeonfish 1.2 & Brown Surgeonfish 6.0; Moorish Idol 1.2 \\
 & Other fish labelled as it (FP) & Brown Surgeonfish 1.3; Acanthuridae 1.2 & Brown Surgeonfish 1.8; Acanthuridae 1.2 \\
Convict Tang & Labelled as (FN) & Ornate Butterfly Fish 1.5; Yellow Tang 1.2 & Ornate Butterfly Fish 0.8; Yellow Tang 0.7 \\
 & Other fish labelled as it (FP) & Ornate Butterfly Fish 1.3; Indo-Pacific Sergeant 1.0 & Ornate Butterfly Fish 5.2; Multiband Butterflyfish 1.2 \\
\bottomrule\end{tabularx}\end{table}

\begin{figure}[H]\centering\includegraphics[width=\linewidth]{figs/fig_err_confidence.pdf}
\caption{Confidence of correct detections (TP), of wrong-species labels given to focus fish (FN), and of focus-species false
positives on other species or on unlabelled regions. All predictions shown have confidence $\geq0.25$ by construction.}\label{fig:conf}\end{figure}

\begin{figure}[H]\centering\includegraphics[width=\linewidth]{figs/fig_err_size.pdf}
\caption{Recall by fish size (shorter box side), baseline and copy-paste; $n$ = labelled fish per bin. Bins with few fish (e.g.\
Convict Tang below 24 px) are not interpretable.}\label{fig:size}\end{figure}

\begin{figure}[H]\centering\includegraphics[width=\linewidth]{figs/fig_err_encounters.pdf}
\caption{Recall per encounter (runs of frames $\le$15 apart), baseline to copy-paste; marker area $\propto$ fish in the encounter.}\label{fig:enc}\end{figure}

\subsection*{Interpretation}
\begin{itemize}
\item \textbf{Baseline failures.} For all three species the largest failure category is a fish with no prediction at confidence
$\geq0.25$ (Forcepsfish 40\%, Orangespine 67\%, Convict 52\% of fish). Wrong-species labels are the second category (25\%, 15\%,
12\%); poor localisation is small (8\%, 3\%, 0\%). More than half of the missed fish (56--64\%) have a correct-species box below the threshold, but
with very low confidence (median 0.02).
\item \textbf{How copy-paste helped differed by species.} Forcepsfish gained mainly through fewer wrong-species labels ($-7.5$\,pp)
and fewer poorly localised boxes ($-2.6$\,pp); the missed share did not change significantly. Convict Tang gained mainly through fewer
missed fish ($-19.2$\,pp), with a smaller drop in wrong-species labels ($-4.0$\,pp). Orangespine Unicornfish changed little
(recall $+4.1$\,pp; no category changed significantly).
\item \textbf{Confusions.} Most Forcepsfish classification errors are Yellow Tang (22.7 of 93 fish per model at baseline, 15.6 with
copy-paste). Orangespine Unicornfish is mostly confused with Brown Surgeonfish, and Convict Tang with Ornate Butterflyfish and
Yellow Tang.
\item \textbf{Cost in false positives.} Copy-paste raised focus-species false positives (pooled 47 $\rightarrow$ 63 per model), mainly
by labelling other fish as Forcepsfish or Convict Tang (e.g.\ Ornate Butterflyfish labelled Convict Tang: 1.3 $\rightarrow$ 5.2 per
model) and by more predictions on unlabelled regions. Precision did not change significantly for any species (Convict Tang
$-5.7$\,pp, interval $-13.1$ to $+7.1$).
\item \textbf{Confidence.} Correct detections became more confident with copy-paste (median Convict Tang 0.74 $\rightarrow$ 0.90,
Forcepsfish 0.76 $\rightarrow$ 0.81). Wrong-species labels and false positives have overlapping confidence distributions, so no single
threshold separates them.
\item \textbf{Size and encounters.} Recall increases with fish size for every species and arm. For Forcepsfish below 16\,px recall is
0.03 (baseline) and 0.05 (copy-paste); the gains are concentrated in fish of 24\,px and larger. Copy-paste improved recall in 10 of
12 Forcepsfish encounters (none worse), all 6 Convict Tang encounters, and 5 of 9 Orangespine encounters (2 worse). One
Forcepsfish encounter (15 fish, median 12\,px) stays near zero under both arms.
\end{itemize}

\subsection*{Supported by the data versus hypotheses}
\textbf{Supported:} the category shares and their changes in Table~\ref{tab:err} (within this footage, these partitions and this
operating point); the dominant confusion pairs; that gains are concentrated in larger fish and spread over most encounters of
Forcepsfish and Convict Tang; that false positives increased while precision did not change significantly.
\textbf{Hypotheses only:} \emph{why} Forcepsfish is confused with Yellow Tang (e.g.\ shared colour) and \emph{why} copy-paste reduces
that confusion (e.g.\ more varied examples) were not tested; whether predictions on unlabelled regions are real fish; whether tiny
fish are missed because of image resolution. These are correlational observations, not causal mechanisms. Results depend on the
chosen thresholds; a different $\tau$ moves fish between ``missed'' and ``correct''.

\section{Does copy-paste help segmentation more than detection?}
The pre-registered tests used box AP50 for detection and mask AP50 for segmentation, so they cannot show that the improvement
magnitudes differ. Here the copy-paste effect ($\Delta$ = copy-paste $-$ baseline, seed-averaged) is computed for every model on
\emph{both} box and mask metrics, and the \emph{difference in effects} between tasks at matched size is bootstrapped with the same
block resamples for both tasks (difference-in-differences).

\begin{table}[H]\centering\caption{Copy-paste effect on focus-species mean AP, by task and size, and the difference in effects
(segmentation box $-$ detection box; $p$ in brackets). Last column: within segmentation, mask effect $-$ box effect.}\label{tab:task}
\footnotesize\setlength{\tabcolsep}{3pt}
\resizebox{\linewidth}{!}{\begin{tabular}{@{}llccccc@{}}\toprule
Metric & Size & Det.\ (box) & Seg.\ (box) & Seg.\ (mask) & Seg.\ box $-$ det.\ box & Seg.\ mask $-$ box\\\midrule
AP50 & Small & \textbf{$+$0.090 [$+$0.026, $+$0.150]} & \textbf{$+$0.112 [$+$0.040, $+$0.176]} & \textbf{$+$0.116 [$+$0.049, $+$0.176]} & $+$0.022 [$-$0.041, $+$0.085] (0.491) & $+$0.005 [$-$0.004, $+$0.014] \\
AP50 & Medium & \textbf{$+$0.089 [$+$0.041, $+$0.144]} & \textbf{$+$0.148 [$+$0.095, $+$0.218]} & \textbf{$+$0.149 [$+$0.099, $+$0.217]} & \textbf{$+$0.060 [$+$0.013, $+$0.110]} (0.026) & $+$0.001 [$-$0.004, $+$0.007] \\
AP50 & Large & \textbf{$+$0.068 [$+$0.015, $+$0.138]} & \textbf{$+$0.089 [$+$0.026, $+$0.153]} & \textbf{$+$0.084 [$+$0.020, $+$0.148]} & $+$0.022 [$-$0.063, $+$0.091] (0.542) & $-$0.006 [$-$0.013, $+$0.000] \\
AP50 & Mean of sizes & -- & -- & -- & $+$0.034 [$-$0.007, $+$0.072] (0.100) & $-$0.000 [$-$0.005, $+$0.004] \\
\midrule
AP50-95 & Small & \textbf{$+$0.071 [$+$0.022, $+$0.124]} & \textbf{$+$0.094 [$+$0.044, $+$0.144]} & \textbf{$+$0.087 [$+$0.036, $+$0.138]} & $+$0.023 [$-$0.025, $+$0.070] (0.326) & $-$0.007 [$-$0.015, $+$0.002] \\
AP50-95 & Medium & \textbf{$+$0.079 [$+$0.037, $+$0.123]} & \textbf{$+$0.121 [$+$0.076, $+$0.181]} & \textbf{$+$0.117 [$+$0.077, $+$0.169]} & \textbf{$+$0.042 [$+$0.005, $+$0.085]} (0.032) & $-$0.004 [$-$0.015, $+$0.007] \\
AP50-95 & Large & \textbf{$+$0.049 [$+$0.006, $+$0.099]} & \textbf{$+$0.085 [$+$0.030, $+$0.140]} & \textbf{$+$0.073 [$+$0.021, $+$0.125]} & $+$0.036 [$-$0.021, $+$0.094] (0.238) & $-$0.012 [$-$0.025, $+$0.002] \\
AP50-95 & Mean of sizes & -- & -- & -- & \textbf{$+$0.034 [$+$0.000, $+$0.066]} (0.049) & \textbf{$-$0.007 [$-$0.013, $-$0.001]} \\
\bottomrule\end{tabular}}\end{table}

\begin{table}[H]\centering\caption{Difference in copy-paste effect per species, segmentation box AP50 $-$ detection box AP50.}\label{tab:taskspecies}
\footnotesize\resizebox{\linewidth}{!}{\begin{tabular}{@{}lcccc@{}}\toprule Species & Small & Medium & Large & Mean of sizes\\\midrule
Forcepsfish & $+$0.053 [$-$0.029, $+$0.125] & $+$0.068 [$-$0.007, $+$0.149] & $-$0.015 [$-$0.112, $+$0.088] & $+$0.035 [$-$0.018, $+$0.086] \\
Orangespine Unicornfish & $-$0.047 [$-$0.129, $+$0.088] & $+$0.015 [$-$0.059, $+$0.135] & $+$0.072 [$-$0.100, $+$0.168] & $+$0.013 [$-$0.035, $+$0.064] \\
Convict Tang & $+$0.059 [$-$0.118, $+$0.217] & \textbf{$+$0.096 [$+$0.012, $+$0.180]} & $+$0.009 [$-$0.124, $+$0.162] & $+$0.054 [$-$0.041, $+$0.146] \\
\bottomrule\end{tabular}}\end{table}

\begin{table}[H]\centering\caption{Decomposition of the copy-paste effect for the pooled focus fish (confidence $\geq0.25$).
Localisation: fish with an any-species box at IoU $\geq0.5$. Identification: correct species among localised fish. Cells:
detection effect / segmentation effect, then segmentation $-$ detection with 95\% interval.}\label{tab:comp}\footnotesize
\begin{tabularx}{\linewidth}{@{}p{4cm}YYY@{}}\toprule Component & Small & Medium & Large\\\midrule
Recall (pp) & $+$7.9 / $+$10.6 \newline \footnotesize{$+$2.7 [$-$3.5, $+$9.7]} & $+$11.3 / $+$15.5 \newline \footnotesize{$+$4.2 [$-$2.3, $+$11.5]} & $+$11.3 / $+$12.6 \newline \footnotesize{$+$1.2 [$-$9.0, $+$9.6]} \\
Localisation (pp) & $+$10.6 / $+$2.5 \newline \footnotesize{\textbf{$-$8.1 [$-$16.4, $-$0.0]}} & $+$5.9 / $+$11.3 \newline \footnotesize{$+$5.4 [$-$1.7, $+$12.6]} & $+$1.2 / $+$10.3 \newline \footnotesize{$+$9.1 [$-$0.3, $+$16.2]} \\
Identification (pp) & $+$4.1 / $+$21.5 \newline \footnotesize{\textbf{$+$17.4 [$+$7.5, $+$28.8]}} & $+$15.6 / $+$16.8 \newline \footnotesize{$+$1.1 [$-$9.2, $+$12.2]} & $+$21.5 / $+$11.1 \newline \footnotesize{\textbf{$-$10.5 [$-$18.8, $-$2.4]}} \\
Mask IoU of correct detections (seg only) & 0.844 $\rightarrow$ 0.854 \newline \footnotesize{$+$0.010 [$-$0.029, $+$0.038]} & 0.858 $\rightarrow$ 0.861 \newline \footnotesize{$+$0.003 [$-$0.027, $+$0.023]} & 0.864 $\rightarrow$ 0.860 \newline \footnotesize{$-$0.003 [$-$0.022, $+$0.012]} \\
\bottomrule\end{tabularx}\end{table}

\begin{figure}[H]\centering\includegraphics[width=\linewidth]{figs/fig_task_effects.pdf}
\caption{Top: copy-paste effect (AP50) for detection (box) and segmentation (box and mask) at each size. Bottom: difference in
effect, segmentation $-$ detection. Filled markers: interval excludes zero.}\label{fig:task}\end{figure}

\begin{figure}[H]\centering\includegraphics[width=\linewidth]{figs/fig_task_components.pdf}
\caption{Copy-paste effect on recall, localisation and identification of focus fish for each set-up.}\label{fig:comp}\end{figure}

\subsection*{Interpretation}
\begin{itemize}
\item \textbf{Metric choice does not explain the apparent difference.} Within segmentation, the mask-AP50 effect equals the box-AP50
effect (mean difference $-0.000$, interval $-0.005$ to $+0.004$). Using mask AP for segmentation therefore did not inflate its
effect at IoU 0.5; at AP50-95 the mask effect is slightly \emph{smaller} than the box effect ($-0.007$).
\item \textbf{The task difference is not robust.} On the same metric (box AP50), the segmentation effect exceeds the detection
effect by $+0.022$ (small), $+0.060$ (medium) and $+0.022$ (large). Only the medium difference has an interval excluding zero
($+0.013$ to $+0.110$, $p=0.026$), which would not survive a correction for three sizes (Bonferroni $p\approx0.08$); averaged over sizes
the difference is $+0.034$ ($-0.007$ to $+0.072$, $p=0.10$). No single species shows a consistent difference.
\item \textbf{Components.} The routes by which copy-paste raised recall differ between set-ups without a consistent task pattern:
at small size segmentation gained more in identification but less in localisation than detection; at large size the opposite;
at medium neither differed. Net recall differences are not significant at any size.
\item \textbf{Mask quality.} Copy-paste did not measurably change mask quality: the mask IoU of correctly detected focus fish was
unchanged at every size (changes $-0.003$ to $+0.010$, all intervals include zero), and every box-matched segmentation detection had
mask IoU $\geq 0.60$ (median 0.88). Mask quality is not the limiting factor for these species.
\end{itemize}
\textbf{Conclusion.} The evidence does not support a task-dependent augmentation effect. Both tasks benefit, by similar amounts once
the same metric is used; the larger numbers for segmentation in the primary table come mainly from one size (medium) and lie within
the statistical uncertainty. Segmentation should not be described as benefiting more. Caveats: detection and segmentation are
different networks trained on different label formats, two seeds only, and seed variance is not in the intervals.

\section{Targeted augmentation versus model capacity}
\begin{table}[H]\centering\caption{Engineering trade-off. Parameters and GFLOPs from one trained checkpoint per architecture (58-class
head), GFLOPs at the 1280-px input used in training and testing. Checkpoint: \texttt{best.pt} (FP16, optimiser stripped). Inference:
median per-image GPU time during the batch-8 test validation on an NVIDIA H100 (Ultralytics ``inference'' timer, excluding pre- and
post-processing) -- a throughput-style measurement on a data-centre GPU, \emph{not} a latency or edge-device benchmark. Training:
mean wall-clock training loop per run; GPU-h: sum of PBS wall time over the 10 runs (one H100 each). AP50: box, pooled CV, mean of
2 seeds.}\label{tab:cost}\footnotesize\setlength{\tabcolsep}{3pt}
\resizebox{\linewidth}{!}{\begin{tabular}{@{}llrrrrrrccc@{}}\toprule
Model & Arm & Params (M) & GFLOPs & Ckpt (MB) & H100 ms/img & Train min/run & GPU-h (10 runs) & Focus AP50 & Non-target AP50$^a$ & Any-fish AP50\\\midrule
Det S & Base & 10.0 & 95 & 20 & 1.9 & 19 & 3.5 & 0.267 & 0.262 & 0.670 \\
 & CP & 10.0 & 95 & 20 & 1.9 & 24 & 4.4 & 0.357 & 0.257 & 0.667 \\
\addlinespace[2pt]
Det M & Base & 21.9 & 305 & 44 & 3.2 & 24 & 4.4 & 0.335 & 0.296 & 0.694 \\
 & CP & 21.9 & 305 & 44 & 3.2 & 40 & 7.1 & 0.424 & 0.295 & 0.687 \\
\addlinespace[2pt]
Det L & Base & 26.3 & 380 & 53 & 4.2 & 38 & 6.5 & 0.373 & 0.294 & 0.681 \\
 & CP & 26.3 & 380 & 53 & 4.1 & 57 & 9.7 & 0.440 & 0.310 & 0.672 \\
\addlinespace[2pt]
Seg S & Base & 11.5 & 153 & 24 & 2.4 & 25 & 4.4 & 0.233 & 0.256 & 0.668 \\
 & CP & 11.5 & 153 & 24 & 2.3 & 40 & 6.8 & 0.345 & 0.260 & 0.658 \\
\addlinespace[2pt]
Seg M & Base & 27.1 & 533 & 55 & 4.1 & 41 & 7.0 & 0.303 & 0.274 & 0.683 \\
 & CP & 27.1 & 533 & 55 & 4.1 & 50 & 8.4 & 0.451 & 0.299 & 0.693 \\
\addlinespace[2pt]
Seg L & Base & 31.5 & 609 & 64 & 4.9 & 53 & 9.2 & 0.351 & 0.290 & 0.671 \\
 & CP & 31.5 & 609 & 64 & 4.9 & 76 & 13.1 & 0.441 & 0.308 & 0.677 \\
\bottomrule\multicolumn{11}{l}{\footnotesize $^a$ Mean AP50 over the 17 well-labelled (${\geq}20$ instances) classes that were not augmented.}
\end{tabular}}\end{table}

\begin{table}[H]\centering\caption{Paired comparisons (box AP50, block bootstrap; bold: interval excludes zero).}\label{tab:cmp}
\footnotesize\setlength{\tabcolsep}{3pt}\resizebox{\linewidth}{!}{\begin{tabular}{@{}llccc@{}}\toprule
Task & Comparison & Focus-species mean & Non-target (17 classes) & Any fish\\\midrule
Detection & CP small $-$ baseline medium & $+$0.022 [$-$0.040, $+$0.088] & \textbf{$-$0.039 [$-$0.077, $-$0.001]} & \textbf{$-$0.026 [$-$0.039, $-$0.012]} \\
Detection & CP medium $-$ baseline large & \textbf{$+$0.051 [$+$0.011, $+$0.107]} & $+$0.001 [$-$0.014, $+$0.029] & $+$0.006 [$-$0.013, $+$0.023] \\
Detection & Augmentation gain (small) & \textbf{$+$0.090 [$+$0.026, $+$0.150]} & $-$0.005 [$-$0.029, $+$0.031] & $-$0.003 [$-$0.016, $+$0.012] \\
Detection & Capacity gain small$\rightarrow$medium & \textbf{$+$0.068 [$+$0.014, $+$0.114]} & \textbf{$+$0.034 [$+$0.014, $+$0.063]} & \textbf{$+$0.023 [$+$0.013, $+$0.035]} \\
Detection & Augmentation gain (medium) & \textbf{$+$0.089 [$+$0.041, $+$0.144]} & $-$0.001 [$-$0.016, $+$0.021] & $-$0.007 [$-$0.019, $+$0.005] \\
Detection & Capacity gain medium$\rightarrow$large & $+$0.037 [$-$0.015, $+$0.090] & $-$0.002 [$-$0.027, $+$0.013] & $-$0.013 [$-$0.025, $+$0.000] \\
\midrule
Segmentation & CP small $-$ baseline medium & $+$0.042 [$-$0.007, $+$0.103] & $-$0.014 [$-$0.033, $+$0.003] & \textbf{$-$0.025 [$-$0.044, $-$0.005]} \\
Segmentation & CP medium $-$ baseline large & \textbf{$+$0.100 [$+$0.048, $+$0.174]} & $+$0.009 [$-$0.015, $+$0.039] & \textbf{$+$0.022 [$+$0.007, $+$0.033]} \\
Segmentation & Augmentation gain (small) & \textbf{$+$0.112 [$+$0.040, $+$0.176]} & $+$0.004 [$-$0.011, $+$0.021] & $-$0.010 [$-$0.023, $+$0.004] \\
Segmentation & Capacity gain small$\rightarrow$medium & \textbf{$+$0.070 [$+$0.011, $+$0.109]} & \textbf{$+$0.019 [$+$0.002, $+$0.038]} & $+$0.015 [$-$0.002, $+$0.033] \\
Segmentation & Augmentation gain (medium) & \textbf{$+$0.148 [$+$0.095, $+$0.218]} & \textbf{$+$0.024 [$+$0.006, $+$0.049]} & $+$0.010 [$-$0.007, $+$0.026] \\
Segmentation & Capacity gain medium$\rightarrow$large & $+$0.048 [$-$0.009, $+$0.104] & $+$0.016 [$-$0.004, $+$0.035] & $-$0.012 [$-$0.026, $+$0.005] \\
\bottomrule\end{tabular}}\end{table}

\begin{figure}[H]\centering\includegraphics[width=\linewidth]{figs/fig_tradeoff.pdf}
\caption{Performance against inference compute (GFLOPs per 1280-px image, log scale) for the baseline and copy-paste version of each
model; arrows show the copy-paste change. Augmentation does not change inference compute.}\label{fig:trade}\end{figure}

\subsection*{Interpretation}
\begin{itemize}
\item \textbf{Inference cost.} Copy-paste leaves architecture, parameters, FLOPs, checkpoint size and measured H100 inference time
unchanged; moving up one size increases GFLOPs by 1.1--3.5$\times$ (e.g.\ detection 95 $\rightarrow$ 305 $\rightarrow$ 380).
\item \textbf{Training cost.} Copy-paste adds 450 synthetic images to about 820 real ones, raising time per epoch by about 50\% and
training time per run by 22--67\% (runs also stop at different epochs) (e.g.\ medium detection 24 $\rightarrow$ 40 min). Generating the synthetic images took 2--5
minutes of CPU time. The whole 120-run experiment used 84.6 H100 GPU-hours.
\item \textbf{Augmented small vs baseline medium.} For the focus species the augmented small models are statistically
indistinguishable from the baseline medium models (detection $+0.022$, $-0.040$ to $+0.088$; segmentation $+0.042$, $-0.007$ to $+0.103$),
but they are worse on non-target species (detection $-0.039$) and on finding fish at all (any-fish $-0.026$ detection, $-0.025$
segmentation). The augmented small model is not a substitute for the medium model.
\item \textbf{Augmented medium vs baseline large.} The augmented medium models outperform the baseline large models on the focus
species (detection $+0.051$, $+0.011$ to $+0.107$; segmentation $+0.100$, $+0.048$ to $+0.174$), with no loss on non-target species or
any-fish AP50, at 80--88\% of the large model's FLOPs.
\item \textbf{Augmentation versus capacity for under-represented species.} At medium size the augmentation gain exceeds the
medium-to-large capacity gain (detection by $+0.051$, segmentation by $+0.100$, both intervals exclude zero); at small size the two
gains are not distinguishable. Capacity, unlike augmentation, also improves the non-target species (small $\rightarrow$ medium:
$+0.034$ detection, $+0.019$ segmentation), whereas augmentation's effect on non-target species is near zero.
\item \textbf{No harm to non-target species.} Augmentation did not significantly reduce non-target AP50 or any-fish AP50 in any of the
six set-ups (largest decreases $-0.005$ and $-0.010$, intervals include zero).
\end{itemize}
\textbf{Conclusion.} For a deployment that must keep a given model size, targeted augmentation improves the under-represented species
at no inference cost and a modest one-off training cost, and at medium size it is a better route to those species than moving to
the large model. It is not a replacement for adding capacity when overall species identification or fish detection matters, and none
of these numbers describe performance on an edge device, which was not benchmarked.

\section{Implications for the report}
\begin{itemize}
\item Replace ``copy-paste gave a gain similar to moving up one model size'' with the paired results: comparable to (not
significantly different from) small$\rightarrow$medium for the target species only, and significantly better than medium$\rightarrow$large.
The current sentence ``small detection with copy-paste (0.357) exceeded medium detection without it (0.335)'' compares point
estimates; the difference is not significant.
\item Do not state that segmentation benefits more; state that both tasks benefit by similar amounts on the same metric.
\item Revise ``most missed Forcepsfish are predicted as background'': at $\tau=0.25$, 55\% of the Forcepsfish not correctly detected at
baseline had no prediction, 34\% were labelled another species (almost all Yellow Tang) and 11\% were poorly localised.
\item TODOs that can now be filled: optimiser AdamW (Ultralytics ``auto'', lr 0.000161) in all 120 runs; AutoBatch sizes 14/8/6
(detection s/m/l) and 12/6/5 (segmentation); Ultralytics 8.4.126, PyTorch 2.13.0; 84.6 H100 GPU-hours (energy upper bound about
59\,kWh at the 700\,W board power, not measured).
\end{itemize}
\end{document}
